\chapter{Stephen Smale --- Fields Medal 1966}\label{ch1966stephensmale:stephen-smale-fields-medal-1966-a-comprehensive-report}

\begin{center}\rule{0.5\linewidth}{0.5pt}\end{center}

\textbf{Prepared:} July 2026 \textbf{Subject:} Stephen Smale
\textbf{Born:} July 15, 1930, Flint, Michigan, USA \textbf{Position:}
Professor, University of California, Berkeley \textbf{Doctoral Advisor:}
Raoul Bott (University of Michigan)

\textbf{Fields Medal Citation:} \emph{``Worked on differential topology
and the theory of dynamical systems; contributed important advances to
the theory of the classification of differentiable manifolds higher
dimensional and the problem of the stability of differentiable mappings.
Proved Poincaré's conjecture for dimensions n ≥ 5.''}

\begin{center}\rule{0.5\linewidth}{0.5pt}\end{center}

\section{1. Introduction}\label{ch1966stephensmale:introduction}

Stephen Smale was awarded the Fields Medal in 1966 for his substantial
contributions to differential topology and dynamical systems.

His proof of the Poincaré conjecture for dimensions \(n \geq 5\) settled
one of the well-known problems in topology and inaugurated modern
manifold classification.

His development of the h-cobordism theorem provided the essential tool
for high-dimensional topology.

In dynamics, Smale's discovery of the Smale horseshoe revealed the
existence of chaotic behavior in deterministic dynamical systems,
founding the modern theory of dynamical systems.

Smale's major contributions include: - The h-cobordism theorem. - The
Poincaré conjecture for \(n \geq 5\). - The Smale horseshoe and chaotic
dynamics. - \(\Omega\)-stability and Axiom A systems. - The Morse--Smale
complex and Morse homology. - Smale's paradox (sphere eversion). -
Smale's 18 problems for the 21st century. - Blum--Shub--Smale
computation theory.

\begin{center}\rule{0.5\linewidth}{0.5pt}\end{center}

\section{2. Biographical Background}\label{ch1966stephensmale:biographical-background}

\subsection{2.1 Early Life and
Education}\label{ch1966stephensmale:early-life-and-education}

Stephen Smale was born on July 15, 1930, in Flint, Michigan.

He entered the University of Michigan in 1948, studying mathematics and
physics.

He completed his Ph.D.~in 1957 under Raoul Bott, writing a thesis on the
topology of regular curves.

\subsection{2.2 Academic Career}\label{ch1966stephensmale:academic-career}

\begin{itemize}
\tightlist
\item
 \textbf{1957:} Ph.D., University of Michigan, under Raoul Bott.
\item
 \textbf{1957--1958:} Instructor, University of Chicago.
\item
 \textbf{1958--1960:} Instructor, Institute for Advanced Study,
 Princeton.
\item
 \textbf{1960--1961:} Professor, UC Berkeley.
\item
 \textbf{1961--1964:} Professor, Columbia University.
\item
 \textbf{1964--1995:} Professor, UC Berkeley.
\item
 \textbf{1966:} Fields Medal at the ICM in Moscow.
\item
 \textbf{1995:} Professor Emeritus, UC Berkeley.
\item
 \textbf{1998:} Wolf Prize in Mathematics.
\end{itemize}

\subsection{2.3 Doctoral Work}\label{ch1966stephensmale:doctoral-work}

Smale's Ph.D.~thesis on regular curves in Riemannian manifolds studied
the topology of the space of immersions.

He proved that a sphere can be turned inside out, now known as Smale's
paradox.

The result was initially met with disbelief because it contradicted
geometric intuition.

\subsection{2.4 Early Influences}\label{ch1966stephensmale:early-influences}

Smale was influenced by the work of Marston Morse and Hassler Whitney.

Their work on critical point theory and immersion theory shaped his
approach to topology.

Raoul Bott's guidance led Smale to apply Morse theory to global problems
in topology.

\begin{center}\rule{0.5\linewidth}{0.5pt}\end{center}

\section{3. The h-Cobordism Theorem}\label{ch1966stephensmale:the-h-cobordism-theorem}

\subsection{3.1 Background}\label{ch1966stephensmale:background}

A \textbf{cobordism} between two \(n\)-manifolds \(M_0\) and \(M_1\) is
a compact \((n+1)\)-manifold \(W\) with boundary
\(\partial W = M_0 \cup M_1\).

It is an \textbf{h-cobordism} if the inclusions
\(M_0 \hookrightarrow W\) and \(M_1 \hookrightarrow W\) are homotopy
equivalences.

\subsection{3.2 Statement of the
Theorem}\label{ch1966stephensmale:statement-of-the-theorem}

The \textbf{h-cobordism theorem} (Smale, 1962) states:

If \(W\) is a simply connected h-cobordism between \(n\)-dimensional
manifolds with \(n \geq 5\), then \(W \cong M_0 \times [0,1]\).

Consequently, \(M_0\) and \(M_1\) are diffeomorphic.

\subsection{3.3 Proof Outline}\label{ch1966stephensmale:proof-outline}

The proof uses \textbf{Morse theory}.

Step 1: Choose a Morse function \(f: W \to [0,1]\) with
\(f^{-1}(0) = M_0\) and \(f^{-1}(1) = M_1\).

Step 2: Cancel critical points in pairs using the \textbf{Whitney trick}
for removing intersections of attaching spheres.

Step 3: Because \(\pi_1 = 0\) and \(n \geq 5\), the Whitney trick works.

Step 4: After canceling all critical points, \(f\) has no critical
points, so \(W\) is a product.

\subsection{3.4 The Whitney Trick}\label{ch1966stephensmale:the-whitney-trick}

The \textbf{Whitney trick} (Whitney, 1944) is a method for removing
unnecessary intersections of submanifolds.

In dimension \(n \geq 5\), two embedded disks of dimensions \(p\) and
\(q\) with \(p+q < n\) can be moved to avoid intersections.

This is the crucial geometric input that fails in low dimensions.

\subsection{3.5 Impact of the h-Cobordism
Theorem}\label{ch1966stephensmale:impact-of-the-h-cobordism-theorem}

The h-cobordism theorem is the central result of high-dimensional
manifold topology.

Together with the Kervaire--Milnor classification of exotic spheres, it
provides a complete classification of simply connected manifolds in
dimensions \(\geq 5\) up to: - Homotopy type (determined by homology). -
The action of the group of homotopy spheres.

\begin{center}\rule{0.5\linewidth}{0.5pt}\end{center}

\section{\texorpdfstring{4. The Poincaré Conjecture for
\(n \geq 5\)}{4. The Poincaré Conjecture for n \textbackslash geq 5}}\label{ch1966stephensmale:the-poincaruxe9-conjecture-for-n-geq-5}

\subsection{4.1 Statement of the
Problem}\label{ch1966stephensmale:statement-of-the-problem}

The \textbf{Poincaré conjecture} states that every simply connected
closed 3-manifold is homeomorphic to \(S^3\).

Smale proved the generalized version: every homotopy \(n\)-sphere
(\(n \geq 5\)) is homeomorphic to \(S^n\).

\subsection{4.2 Smale's Proof}\label{ch1966stephensmale:smales-proof}

The proof follows directly from the h-cobordism theorem.

If \(\Sigma^n\) is a homotopy \(n\)-sphere, remove two small balls from
\(\Sigma^n\).

The resulting manifold \(W\) is an h-cobordism between \(S^{n-1}\) and
\(S^{n-1}\).

By the h-cobordism theorem, \(W \cong S^{n-1} \times [0,1]\).

Glue the balls back in: \(\Sigma^n\) is a union of two \(n\)-disks along
their boundaries, hence \(S^n\).

\subsection{4.3 Impact of the Proof}\label{ch1966stephensmale:impact-of-the-proof}

This work revolutionized topology.

It showed that high dimensions (\(n \geq 5\)) are more tractable than
low dimensions, where the Whitney trick fails.

The Poincaré conjecture in dimension 3 was finally proved by Perelman in
2003 (Fields Medal 2006), using Ricci flow.

The \(n=4\) case was proved by Freedman (Fields Medal 1986) using Casson
handles.

\subsection{4.4 Historical Context}\label{ch1966stephensmale:historical-context}

The Poincaré conjecture was one of the well-known problems in all of
mathematics.

Smale's work in dimension \(\geq 5\) opened the door to understanding
why different dimensions behave so differently in manifold topology.

\begin{center}\rule{0.5\linewidth}{0.5pt}\end{center}

\section{5. The Smale Horseshoe}\label{ch1966stephensmale:the-smale-horseshoe}

\subsection{5.1 Definition}\label{ch1966stephensmale:definition}

The \textbf{Smale horseshoe} is a map \(f: D \to \mathbb{R}^2\) defined
on a square \(D\).

The map contracts vertically, stretches horizontally, and folds the
square into a horseshoe shape.

The invariant set

\[
\Lambda = \bigcap_{n \in \mathbb{Z}} f^n(D)
\]

is a Cantor set, and the dynamics on \(\Lambda\) is topologically
conjugate to the \textbf{shift map} on two symbols \(\Sigma_2\).

\subsection{5.2 Chaotic Dynamics}\label{ch1966stephensmale:chaotic-dynamics}

The horseshoe exhibits: - \textbf{Sensitive dependence on initial
conditions}: nearby points diverge exponentially. - \textbf{A dense
orbit}: there is a point whose forward orbit is dense in \(\Lambda\). -
\textbf{Periodic orbits are dense}: periodic points are dense in
\(\Lambda\). - \textbf{Topological entropy}:

\[
h_{\text{top}}(f|_\Lambda) = \log 2.
\]

These are the defining characteristics of chaos.

\subsection{5.3 Symbolic Dynamics}\label{ch1966stephensmale:symbolic-dynamics}

The shift map on two symbols is defined as:

\[
\sigma: \Sigma_2 \to \Sigma_2,\quad
\sigma(s_0 s_1 s_2 \dots) = s_1 s_2 s_3 \dots
\]

Each symbol sequence codes a unique orbit in the horseshoe.

Periodic orbits correspond to periodic sequences.

The coding provides a complete description of the dynamics on
\(\Lambda\).

\subsection{5.4 Impact of the Horseshoe}\label{ch1966stephensmale:impact-of-the-horseshoe}

The horseshoe is the paradigm for chaotic dynamics in deterministic
systems.

It appears in the study of: - The Lorenz attractor and strange
attractors. - The homoclinic tangle and Melnikov methods. -
Period-doubling cascades and universality. - Symbolic dynamics and
coding of orbits.

\begin{center}\rule{0.5\linewidth}{0.5pt}\end{center}

\section{\texorpdfstring{6. Axiom A and
\(\Omega\)-Stability}{6. Axiom A and \textbackslash Omega-Stability}}\label{ch1966stephensmale:axiom-a-and-omega-stability}

\subsection{6.1 Background}\label{ch1966stephensmale:background-1}

A diffeomorphism \(f: M \to M\) is \textbf{structurally stable} if every
sufficiently close diffeomorphism \(g\) is topologically conjugate to
\(f\).

The \textbf{non-wandering set} \(\Omega(f)\) consists of points \(x\)
such that for every neighborhood \(U\) of \(x\), there exists \(n > 0\)
with \(f^n(U) \cap U \neq \varnothing\).

\subsection{6.2 Axiom A}\label{ch1966stephensmale:axiom-a}

Smale introduced \textbf{Axiom A}: 1. The non-wandering set
\(\Omega(f)\) is hyperbolic. 2. Periodic points are dense in
\(\Omega(f)\).

A set is \textbf{hyperbolic} if the tangent bundle restricted to the set
splits into contracting and expanding directions.

\subsection{\texorpdfstring{6.3 The \(\Omega\)-Stability
Theorem}{6.3 The \textbackslash Omega-Stability Theorem}}\label{ch1966stephensmale:the-omega-stability-theorem}

Smale proved: if \(f\) satisfies Axiom A and the \textbf{strong
transversality condition} (stable and unstable manifolds intersect
transversely), then \(f\) is \(\Omega\)-stable.

This means: any nearby \(g\) has \(\Omega(g)\) homeomorphic to
\(\Omega(f)\), and \(f|_{\Omega(f)}\) is topologically conjugate to
\(g|_{\Omega(g)}\).

\subsection{6.4 The Spectral Decomposition
Theorem}\label{ch1966stephensmale:the-spectral-decomposition-theorem}

For an Axiom A diffeomorphism, the non-wandering set decomposes into a
finite disjoint union of closed invariant sets (basic sets), each
topologically transitive.

\subsection{6.5 The Palis--Smale Program}\label{ch1966stephensmale:the-palissmale-program}

Smale outlined a program for understanding the dynamics of typical
diffeomorphisms.

The program proposed: 1. Hyperbolic systems are structurally stable. 2.
Generic systems satisfy Axiom A. 3. The non-wandering set decomposes
into basic sets with a dense orbit.

Part (2) turned out to be false (Newhouse showed wild hyperbolic sets
exist), but the program shaped dynamical systems for decades.

\begin{center}\rule{0.5\linewidth}{0.5pt}\end{center}

\section{7. Morse Homology}\label{ch1966stephensmale:morse-homology}

\subsection{7.1 The Morse--Smale Complex}\label{ch1966stephensmale:the-morsesmale-complex}

Smale's work on the h-cobordism theorem led him to study gradient flows
of Morse functions on manifolds.

A Morse function with the \textbf{Smale condition} (gradient flow has no
connections between critical points of the same index) defines a
\textbf{Morse--Smale complex}.

The chain groups are generated by critical points.

The boundary operator counts gradient flow lines.

\subsection{7.2 Morse Homology}\label{ch1966stephensmale:morse-homology-1}

The \textbf{Morse homology} \(H_*(M, \mathbb{Z})\) is isomorphic to the
singular homology of \(M\).

This provided: - A new proof of the Morse inequalities. - A
combinatorial approach to homology. - Connections to Floer homology and
symplectic geometry.

\begin{center}\rule{0.5\linewidth}{0.5pt}\end{center}

\section{8. Smale's Paradox}\label{ch1966stephensmale:smales-paradox}

\subsection{8.1 Statement}\label{ch1966stephensmale:statement}

\textbf{Smale's paradox} (1958) states that the 2-sphere \(S^2\) can be
turned inside out by a regular homotopy of immersions.

There exists a regular homotopy

\[
F: S^2 \times [0,1] \to \mathbb{R}^3
\]

such that \(F_0\) is the standard inclusion and \(F_1\) is the antipodal
map (the sphere turned inside out).

\subsection{8.2 The Proof}\label{ch1966stephensmale:the-proof}

Smale proved that the space of immersions

\[
\text{Imm}(S^2, \mathbb{R}^3)
\]

is connected, using the h-principle and Morse theory.

The result was initially met with disbelief because it contradicts
geometric intuition.

\subsection{8.3 Visualization}\label{ch1966stephensmale:visualization}

The eversion can be visualized through a series of intermediate stages.

The sphere passes through itself (which is allowed for immersions).

A famous video by Nelson Max from the National Center for Supercomputing
Applications visualizes the eversion.

\begin{center}\rule{0.5\linewidth}{0.5pt}\end{center}

\section{9. Complexity and
Computation}\label{ch1966stephensmale:complexity-and-computation}

\subsection{9.1 The Hirsch Conjecture}\label{ch1966stephensmale:the-hirsch-conjecture}

Smale studied the \textbf{Hirsch conjecture}: the diameter of the graph
of a convex polytope is bounded by \(f - d\).

This is a problem in linear programming and combinatorial optimization.

While the conjecture is false in general, Smale's work led to smoothed
analysis of the simplex algorithm.

\subsection{9.2 Smale's 18 Problems}\label{ch1966stephensmale:smales-18-problems}

In 1998, Smale published a list of 18 problems for the 21st century.

The list includes: - The Riemann hypothesis. - The Poincaré conjecture
(since solved). - \(P\) vs \(NP\) (computational complexity). - The
Jacobian conjecture. - The Birch--Swinnerton-Dyer conjecture. - The
complexity of polynomial systems.

\begin{center}\rule{0.5\linewidth}{0.5pt}\end{center}

\section{10. The Blum--Shub--Smale Model}\label{ch1966stephensmale:the-blumshubsmale-model}

\subsection{10.1 Definition}\label{ch1966stephensmale:definition-1}

\textbf{BSS computation theory} (with Blum and Shub) generalizes the
Turing machine model to computation over rings.

Instead of discrete bits, computations operate on real numbers, with
exact arithmetic and branching on equalities.

\subsection{10.2 Main Results}\label{ch1966stephensmale:main-results}

\begin{itemize}
\tightlist
\item
 \(P_{\mathbb{R}}\) vs \(NP_{\mathbb{R}}\): the problem over the reals
 differs substantially from the Turing machine version.
\item
 The \textbf{nullstellensatz} provides algebraic certificates for
 \(NP_{\mathbb{R}}\) problems.
\item
 The complexity of polynomial systems is a central open question.
\end{itemize}

\begin{center}\rule{0.5\linewidth}{0.5pt}\end{center}

\section{11. Applications to Mechanics}\label{ch1966stephensmale:applications-to-mechanics}

\subsection{11.1 The N-Body Problem}\label{ch1966stephensmale:the-n-body-problem}

Smale applied topology to problems in mechanics, particularly the N-body
problem.

He used methods from dynamical systems to study the topology of the
configuration space.

He showed that the topology of the space of N-body configurations
influences possible motions.

\subsection{11.2 Periodic Orbits}\label{ch1966stephensmale:periodic-orbits}

Smale studied the existence of periodic orbits in the N-body problem.

His work connected celestial mechanics to the theory of dynamical
systems, showing that the complexity of the horseshoe appears in
gravitational problems.

\begin{center}\rule{0.5\linewidth}{0.5pt}\end{center}

\section{12. Awards and Recognition}\label{ch1966stephensmale:awards-and-recognition}

\textbf{1966:} - Fields Medal (IMU).

\textbf{1988:} - Chauvenet Prize (MAA).

\textbf{1996:} - National Medal of Science (USA).

\textbf{1998:} - Wolf Prize in Mathematics.

\textbf{2007:} - Lester R. Ford Award.

\textbf{Elected to:} - National Academy of Sciences (USA). - American
Academy of Arts and Sciences. - Royal Society of London, Foreign Member.
- French Academy of Sciences, Foreign Member.

\begin{center}\rule{0.5\linewidth}{0.5pt}\end{center}

\section{13. Selected Major
Publications}\label{ch1966stephensmale:selected-major-publications}

\begin{enumerate}
\def\labelenumi{\arabic{enumi}.}
\item
 \textbf{``Generalized Poincaré's conjecture in dimensions greater than
 four.''} \emph{Annals of Mathematics} 74, 391--406, 1961. ---
 \emph{The proof of the Poincaré conjecture for \(n \geq 5\).}
\item
 \textbf{``On the structure of manifolds.''} \emph{American Journal of
 Mathematics} 84, 387--399, 1962. --- \emph{The h-cobordism theorem.}
\item
 \textbf{``Differentiable dynamical systems.''} \emph{Bulletin of the
 AMS} 73, 747--817, 1967. --- \emph{Foundational survey introducing the
 horseshoe and Axiom A.}
\item
 \textbf{``Morse inequalities for a dynamical system.''} \emph{Bulletin
 of the AMS} 66, 43--49, 1960. --- \emph{The Morse--Smale complex.}
\item
 \textbf{``A classification of immersions of the two-sphere.''}
 \emph{Transactions of the AMS} 90, 281--290, 1959. --- \emph{Smale's
 paradox: sphere eversion.}
\item
 \textbf{``The \(\Omega\)-stability theorem.''} (with J. Palis)
 \emph{Proceedings of Symposia in Pure Math.} 14, AMS, 1970. ---
 \emph{Fundamental theorem on structural stability.}
\item
 \textbf{``On the efficiency of algorithms of analysis.''} (with L.
 Blum and M. Shub) \emph{Bulletin of the AMS} 13, 87--91, 1985. ---
 \emph{BSS computation theory.}
\item
 \textbf{``The horseshoe map.''} In \emph{The Mathematics of Time},
 Springer, 1980. --- \emph{Definitive account of the horseshoe.}
\item
 \textbf{``Topology and mechanics I--II.''} \emph{Inventiones
 Mathematicae} 10, 305--331, 1970; 11, 45--64, 1970. ---
 \emph{Applications of topology to mechanics.}
\item
 \textbf{``Mathematical problems for the next century.''}
 \emph{Mathematical Intelligencer} 20, 7--15, 1998. --- \emph{Smale's
 18 problems.}
\end{enumerate}

\begin{center}\rule{0.5\linewidth}{0.5pt}\end{center}

\section{14. References and Further
Reading}\label{ch1966stephensmale:references-and-further-reading}

\subsection{Official Sources}\label{ch1966stephensmale:official-sources}

\begin{itemize}
\tightlist
\item
 IMU Fields Medal 1966:
 https://www.mathunion.org/imu-awards/fields-medal/fields-medals-1966
\item
 Wolf Prize Citation.
\end{itemize}

\subsection{Key Technical References}\label{ch1966stephensmale:key-technical-references}

\begin{itemize}
\tightlist
\item
 Smale, S. \emph{The Mathematics of Time}. Springer, 1980.
\item
 Smale, S. ``Differentiable dynamical systems.'' \emph{Bulletin AMS}
 73, 1967.
\item
 Milnor, J. \emph{Lectures on the h-Cobordism Theorem}. Princeton,
 1965.
\end{itemize}

\subsection{Textbooks and Surveys}\label{ch1966stephensmale:textbooks-and-surveys}

\begin{itemize}
\tightlist
\item
 Hirsch, M. \emph{Differential Topology}. Springer, 1976.
\item
 Palis, J. and de Melo, W. \emph{Geometric Theory of Dynamical
 Systems}. Springer, 1982.
\item
 Katok, A. and Hasselblatt, B. \emph{Introduction to the Modern Theory
 of Dynamical Systems}. Cambridge, 1995.
\end{itemize}

\subsection{Biography and Exposition}\label{ch1966stephensmale:biography-and-exposition}

\begin{itemize}
\tightlist
\item
 Bott, R. ``The Work of Stephen Smale.'' ICM 1966 Laudatio.
\item
 Hirsch, M. ``Stephen Smale.'' \emph{Notices of the AMS}, 1990.
\item
 Shub, M. ``The Work of Stephen Smale in Dynamics.'' \emph{Bulletin
 AMS}, 1998.
\end{itemize}

\begin{center}\rule{0.5\linewidth}{0.5pt}\end{center}

\emph{(End of report)}
